Las conclusiones siguen el orden de los objetivos y se limitan a lo que los capítulos de Resultados y Discusión sustentan. Cada una indica su alcance.

\textbf{Objetivo general.} El objetivo general se cumplió de forma parcial. La librería existe, intercepta las peticiones de Express, clasifica las cuatro categorías de ataque y bloquea con Random Forest. La detección de variantes desconocidas solo tiene evidencia indirecta: en validación, Isolation Forest marca 19 de los 39 ataques que Random Forest deja pasar, y bloquear sus marcas agregaría 866 benignos bloqueados. Por ese costo, Isolation Forest solo registra. El requisito de no comprometer el rendimiento se cumple para el criterio de latencia en carga normal y de ataque, y no se cumple para memoria, CPU ni carga de volumen.

\textbf{Objetivo específico 1.} El extractor cumple el criterio en aislamiento: el p95 con \texttt{worker\_threads} es de 0.0273~ms, frente al límite de 5~ms. El sobrecosto del middleware completo es de 12.08~ms en carga normal y de 12.85~ms en carga de ataque, y la extracción no lo explica. Sebastián Huertas implementó el extractor inicial de 72 características. El autor analizó y justificó cada característica, definió cuáles no existen en inferencia, fijó el contrato de entrada y añadió cuatro características y las correcciones derivadas del análisis del modelo.

\textbf{Objetivo específico 2.} Ambos criterios se cumplen con margen sobre el conjunto de prueba: el F1 macro de Random Forest es de 0.9776, con un mínimo de 0.9239 en \texttt{cmdi}, y Isolation Forest alcanza un \textit{recall} de 0.9125 con una tasa de falsos positivos de 0.0546. La paridad entre scikit-learn y ONNX es del orden de $10^{-7}$. Estas cifras tienen tres límites: el conjunto de prueba se leyó tres veces para la generación vigente, las métricas se calculan por argmax sin el umbral de 0.35, y las particiones comparten la estructura de casi-duplicados del entrenamiento. Miden, por tanto, la separación entre clases dentro del corpus y no la generalización a tráfico externo. Con \texttt{rf\_v10} e \texttt{if\_v9}, sobre otra partición de prueba, los criterios también se cumplen (F1 macro de 0.9787; anexo~\ref{anexo:v10}).

\textbf{Objetivo específico 3: cobertura.} En el corpus sin solapamiento relevante, la cobertura supera el 80\,\% con ambos pares de modelos: con \texttt{rf\_v11}/\texttt{if\_v10} es de 99.5\,\% en \texttt{sqli}, 95.0\,\% en \texttt{path\_traversal} y 96.0\,\% en \texttt{cmdi}, y con el paquete publicado, de 100\,\%, 85.0\,\% y 91.5\,\%. Con el paquete publicado, el límite inferior del intervalo de \texttt{path\_traversal} (79.4\,\%) queda bajo el criterio. La categoría \texttt{xss} carece de corpus limpio, y su cobertura de 97.3\,\% proviene del corpus de 590 payloads. En defensa en capas, logSguarDian capturó 37 de los 40 ataques que ModSecurity dejó pasar, y la comparación no se repitió con el corpus sin solapamiento.

\textbf{Objetivo específico 3: latencia y recursos.} El criterio de $\Delta$p95 $\leq$ 15~ms se cumple en carga normal (12.08~ms) y de ataque (12.85~ms), y no en volumen, donde el $\Delta$p95 es de 6\,495.46~ms y, en ese escenario con el middleware activo, el 7.00\,\% de las peticiones agota el tiempo de espera. Los datos provienen de una corrida con cinco repeticiones en un ejecutor de 2 vCPU. Los criterios del protocolo de memoria adicional (150~MB) y de CPU (15\,\%) no se cumplen: el incremento de memoria va de 331.9 a 375.9~MiB según el escenario, y el de CPU es de +529.6\,\% a +614.4\,\% en carga normal y de ataque. El costo por componente no se midió.

\textbf{Falsos positivos y generalización.} El corpus benigno es la debilidad principal. Con un User-Agent de navegador, el modelo bloqueó el 97.1\,\% de las peticiones benignas en OWASP Juice Shop (303 de 312) y el 92.5\,\% en \texttt{dvws-node} (37 de 40; el modelo exacto de esa medición no consta). En la aplicación de referencia, con otro modelo, ninguna petición benigna con el User-Agent de Artillery recibió un 403 en los escenarios normal y de volumen. La exclusión por fuente con la configuración de \texttt{rf\_v11} sigue el mismo sentido: los falsos positivos benignos suben a 76.0\,\%, 27.0\,\% y 94.5\,\% al excluir \texttt{modsec\_learn}, \texttt{payloads\_csv} o \texttt{xss\_dataset} del entrenamiento (anexo~\ref{anexo:loso}). El User-Agent es la característica de mayor peso en las clases \texttt{benign} y \texttt{sqli}, y solo el 0.4\,\% de las filas benignas del corpus tiene un User-Agent largo.

\textbf{Alcance de lo entregado.} El paquete publicado (\texttt{logsguardian@0.1.0}) contiene \texttt{rf\_v10} e \texttt{if\_v9}, y las métricas de modelo de este trabajo corresponden a \texttt{rf\_v11} e \texttt{if\_v10}, que no están publicados. El trabajo documenta la procedencia de las cifras principales mediante hashes, locks de las particiones y anexos, y declara las limitaciones y los datos que no constan. Publicar una versión con los modelos evaluados y repetir las mediciones de integración con ella queda pendiente.
